\begin{abstract}
Representation engineering shows that large language models encode sycophancy as a linear direction in activation space. A contrastive sycophancy direction (\dsyc{}) extracted from style-controlled pairs achieves cross-validated area under the ROC curve (AUROC) above 0.82, which suggests it could serve as a reward signal for reinforcement learning against sycophantic behavior. We test this hypothesis with Group Relative Policy Optimization (GRPO) under three \dsyc{}-based reward formulations across two model families. Every formulation fails. The cosine reward cannot separate valid revision from sycophantic capitulation in the narrow online activation gap. The probe reward creates a stubbornness trap: the model learns that never revising its answer maximizes reward, and the true-positive flip rate drops to zero within 500 steps. When we remove \dsyc{} entirely and train with a purely behavioral reward, sycophancy scores plateau below the vanilla baseline. Only supervised fine-tuning with dense per-token supervision succeeds. The root cause is a semantic gap between discriminability and reward utility: \dsyc{} encodes \emph{whether} belief changed, not \emph{whether it should have}. Valid revision and sycophantic capitulation produce overlapping activations because both involve the model shifting its stated position toward the user's. A single scalar reward per completion, whether from representations or behavioral classifiers, cannot convey which belief changes are warranted. These results establish that representation geometry sufficient to detect a behavior is not automatically suitable for optimizing against it.
\end{abstract}
